% Generated draft: requires mathematical and source review.
\documentclass[11pt]{article}
\usepackage[margin=25mm]{geometry}
\usepackage{fontspec}
\setmainfont{Latin Modern Roman}
\usepackage{amsmath,amssymb,mathtools,unicode-math}
\setmathfont{Latin Modern Math}
\usepackage{enumitem,needspace,xurl,hyperref}
\setlength{\parindent}{0pt}
\setlength{\parskip}{5pt}
\setlength{\emergencystretch}{3em}
\newcommand{\N}{\mathbb{N}}
\newcommand{\Z}{\mathbb{Z}}
\newcommand{\Q}{\mathbb{Q}}
\newcommand{\R}{\mathbb{R}}
\newcommand{\C}{\mathbb{C}}
\newcommand{\F}{\mathbb{F}}
\newcommand{\A}{\mathbb{A}}
\newcommand{\PP}{\mathbb{P}}
\newcommand{\E}{\mathbb{E}}
\newcommand{\eps}{\varepsilon}
\newcommand{\dd}{\,\mathrm{d}}
\begin{document}
\section*{UnsolvedMath ID: 30004639}
\textbf{Source problem code:} catalogue entry
\subsection{Definitions and mathematical statement}
The supplied record poses a research programme rather than a fully specified mathematical conjecture. Its principal request is to find suitable admissibility criteria for physically reasonable solutions of multidimensional hyperbolic balance laws, especially the Euler system.

For notation only, a chosen deterministic balance law in conservation form may be written as
\[
\partial_t u+\sum_{j=1}^{d}\partial_{x_j}F_j(u)=S(t,x,u),\qquad (t,x)\in I\times\Omega,
\]
where \(d\geq2\), \(\Omega\subseteq\R^d\) is open, \(I\subseteq\R\) is an open time interval, \(u:I\times\Omega\to U\subseteq\R^m\) is the state, \(F_j:U\to\R^m\) are fluxes, and \(S:I\times\Omega\times U\to\R^m\) is a source term. These objects are placeholders, not model specifications supplied by the record. When \(u\), \(F_j(u)\), and \(S(t,x,u)\) are locally integrable, the distributional balance law means
\[
\int_{I\times\Omega}\left[u\cdot\partial_t\varphi+\sum_{j=1}^{d}F_j(u)\cdot\partial_{x_j}\varphi+S(t,x,u)\cdot\varphi\right]\,\dd t\,\dd x=0
\]
for every \(\varphi\in C_c^\infty(I\times\Omega;\R^m)\). This interior formulation specifies neither initial nor boundary conditions. It also does not prescribe a stochastic interpretation for random systems.

An admissibility criterion would be an additional condition selecting physically reasonable solutions from the solutions of a specified balance law. The supplied research tasks are to identify suitable such conditions, especially for the Euler system; determine their implications for multiphase-flow models, random systems, and structure-preserving numerical schemes; and assess whether hyperbolic balance laws can predict pandemic evolution.

\textbf{Unresolved formulation.} The record specifies no particular fluxes, source terms, state space, hyperbolicity assumptions, Euler variant or constitutive closure, initial or boundary data, solution class, or definition of physical reasonableness. It does not identify the multiphase or random models, numerical structures to preserve, or a pandemic model, prediction target, time horizon, data, or accuracy criterion. Consequently, it does not determine a single quantified existence, uniqueness, stability, convergence, or prediction assertion. In particular, the request must not be strengthened to demand one universal admissibility criterion for all the mentioned models. The displayed equation is explanatory notation only; the unresolved choices require human mathematical review.
\subsection{Short English statement}
Find admissibility conditions selecting physically reasonable solutions of multidimensional hyperbolic balance laws, particularly Euler systems. Study their consequences for multiphase flow, random systems, and structure-preserving numerical methods, and assess their potential for predicting pandemic evolution. The supplied record leaves the models and success criteria unspecified, so this remains a research programme requiring further formulation.
\subsection{Sources}
\begin{itemize}
\item Local UnsolvedMath registry: supplied complete source record, permanent ID \(30004639\), source problem code \texttt{catalogue entry}, title ``Admissibility and Prediction for Multidimensional Hyperbolic Balance Laws,'' in the collection ``Oberwolfach Reports Open Problems.'' No individual proposer is supplied; the proposed year is recorded as \(2021\).
\item Supplied report citation: \emph{Hyperbolic Balance Laws: modeling, analysis, and numerics} (2022), Oberwolfach Reports, pp.\ 589--661, DOI \texttt{10.4171/owr/2021/11}. Supplied source URL: \url{https://doi.org/10.4171/owr/2021/11}. The background attributes the question to workshop open-problem discussions documented in Oberwolfach Reports; no specific contribution, author, or problem page within the report is supplied.
\item Supplied literature source: ``Oberwolfach Report 11/2021,'' \url{https://doi.org/10.4171/owr/2021/11}. The exact supplied literature-check date is \texttt{2026-08-22}. The record labels the problem open and supplies the assessment that no universal admissibility criterion or well-posedness theory is known for multidimensional hyperbolic balance laws and Euler systems across the listed models. This is supplied dated curation, not an independently verified literature assessment. The background explicitly requests independent verification of time-sensitive status and cited claims.
\item The registry marks the statement as ``exact'' and says it was retained byte-for-byte without statement-level reconstruction. These are supplied source-verification claims; neither the report text nor that claimed correspondence has been independently checked here. The citation's year \(2022\), the proposed year \(2021\), and the report designation 11/2021 are preserved as supplied without resolving their relationship.
\end{itemize}
\end{document}
